\section{Mathematical Formalization of the Stable Effective Volume Metric (EV3)}
\label{app:ev3_formal}

This appendix provides a rigorous mathematical derivation of the Stable Effective Volume ($\text{EV3}$) metric, establishes its axiomatic foundations, proves its fundamental invariance and stability properties, and explains why previous determinant-based computations led to numerical underflow.

\subsection{Axiomatic Derivation}
Let $X \in \mathbb{R}^{N \times D}$ represent an activation matrix obtained from a neural layer, consisting of $N$ samples embedded in a $D$-dimensional representation space. We first define the centered activation matrix $\tilde{X}$:
\begin{equation}
\tilde{X} = X - \frac{1}{N} \mathbf{1}_{N \times N} X
\end{equation}
where $\mathbf{1}_{N \times N}$ is the identity-like matrix of ones. The empirical covariance matrix of the activations is given by:
\begin{equation}
C = \frac{1}{N} \tilde{X}^T \tilde{X}
\end{equation}
By performing Singular Value Decomposition (SVD) on the centered activation matrix $\tilde{X}$, we obtain:
\begin{equation}
\tilde{X} = U \Sigma V^T
\end{equation}
where $U \in \mathbb{R}^{N \times N}$ and $V \in \mathbb{R}^{D \times D}$ are orthogonal matrices, and $\Sigma \in \mathbb{R}^{N \times D}$ contains the singular values $\sigma_1 \ge \sigma_2 \ge \dots \ge \sigma_k \ge 0$ on its diagonal, with $k = \min(N, D)$. The eigenvalues $\lambda_i$ of the covariance matrix $C$ are directly related to the singular values by $\lambda_i = \sigma_i^2 / N$.

The principal axes of the representation uncertainty ellipsoid are proportional to the singular values $\sigma_i$. Conceptually, the volume of this ellipsoid is proportional to the product of its semi-axes:
\begin{equation}
\text{Vol}(\mathcal{E}) \propto \prod_{i=1}^k \sigma_i
\end{equation}
To ensure stability across varying dimensions and prevent numerical underflow, we define $\text{EV3}$ as the dimension-normalized geometric mean of the significant singular values:
\begin{equation}
\text{EV3}(X) = \exp \left( \frac{1}{M} \sum_{i=1}^M \log \sigma_i \right)
\end{equation}
where $M \le k$ is the number of active singular values that strictly exceed a regularizing threshold $\epsilon = 10^{-10}$. If all singular values are below the threshold ($M = 0$), we define $\text{EV3}(X) = 0.0$.

\subsection{Mathematical Properties and Proofs}
We prove four fundamental properties of $\text{EV3}$ that validate its use as a representation audit metric:

\begin{enumerate}
    \item \textbf{Rotation Invariance:} Let $R \in \text{SO}(D)$ be an orthogonal rotation matrix. The rotated embeddings are $X' = X R$. Since $R$ is orthogonal ($R^T R = I$), the centered matrix is $\tilde{X}' = \tilde{X} R$. The SVD of the rotated matrix is:
    \begin{equation}
    \tilde{X}' = \tilde{X} R = U \Sigma (V^T R)
    \end{equation}
    Since the product of two orthogonal matrices $V^T R$ is also orthogonal, the singular values $\Sigma$ remain completely unchanged. Consequently, $\text{EV3}(X R) = \text{EV3}(X)$.
    
    \item \textbf{Non-Negativity:} Since each active singular value satisfies $\sigma_i > \epsilon > 0$, their logarithms are real, and their geometric mean $\text{EV3}$ is strictly positive ($\text{EV3} > 0$). When the representation completely collapses to a single point, $M=0$, and $\text{EV3} = 0.0$, satisfying non-negativity.
    
    \item \textbf{Subspace Stability (Avoiding Determinant Collapse):} In deep networks, representations often collapse to a lower-dimensional manifold of dimension $d < D$, meaning the remaining $D - d$ singular values drop to zero. A standard determinant $\det(C) = \prod_{i=1}^D \lambda_i$ collapses immediately to $0$. In contrast, $\text{EV3}$ filters out the collapsed dimensions ($\sigma_i \le \epsilon$) and measures the stable geometric volume spanned strictly in the active subspace of dimension $d$.
    
    \item \textbf{White Noise Limit:} For perfectly uncorrelated, isotropic Gaussian activations with variance $\sigma^2$ in all directions, all active singular values are identical ($\sigma_i \approx \sigma\sqrt{N}$). The log-space mean yields $\text{EV3} \approx \sigma\sqrt{N}$, which behaves linearly with signal scale.
\end{enumerate}

\subsection{Numerical Stability Analysis and Comparison}
In previous implementations, the representation volume was computed directly as the determinant of the covariance matrix, $\det(C)$. If the network had $D=64$ features and the singular values in noise directions were around $0.05$, the product of these eigenvalues resulted in:
\begin{equation}
\det(C) = \prod_{i=1}^{64} \lambda_i \approx 0.05^{64} \approx 10^{-84}
\end{equation}
For larger dimensions ($D \ge 100$), this product rapidly exceeded the double-precision floating-point limit ($10^{-308}$), causing catastrophic numerical underflow to $0.0$ or producing arbitrary values in the range $10^{-213}$ to $10^{-221}$. By evaluating the geometric mean in log-space ($\frac{1}{M}\sum \log\sigma_i$) and then exponentiating, $\text{EV3}$ remains stably bounded in the range $[0.1, 100.0]$ regardless of embedding dimension.

Table~\ref{tab:comparison} compares $\text{EV3}$ with existing representation alignment metrics.

\begin{table}[h]
\centering
\caption{Comparison of representation audit metrics.}
\label{tab:comparison}
\begin{tabular}{lccc}
\toprule
Metric & Evaluates Alignment & Evaluates Vol/Complexity & Numerical Stability \\
\midrule
CKA & Yes & No & High \\
SVCCA & Yes & No & Medium \\
PWCCA & Yes & No & Medium \\
\textbf{EV3 (New)} & \textbf{No} & \textbf{Yes} & \textbf{High} \\
\bottomrule
\end{tabular}
\end{table}
